% Mechanical relocation generated from the preservation ledger.
% Existing prose is included byte-for-byte from preserved blocks.

\section{Grouped Holdout Validation}
\label{sec:grouped-holdout-validation}

Holdout validation asks whether a model fitted without selected observations
predicts those untouched observations. Measurement rows must not be split
independently when they share a vehicle journey, stop time series, or another
operational unit: such a split leaks closely related information between
calibration and validation. The grouping unit must be selected before examining
results and must correspond to the intended generalization question.

For a partition $\mathcal M=\mathcal M_{\mathrm{cal}}\cup
\mathcal M_{\mathrm{hold}}$, estimation uses likelihood contributions only
from $\mathcal M_{\mathrm{cal}}$. Predictions are computed for both sets, but
held-out counts must have no effect on fitted parameters. This invariance is a
required validation check. Calibration and holdout results should report the
same metrics and grouping definitions.

Holding out complete vehicle journeys tests transfer to other journeys in the
represented timetable; holding out stops tests spatial transfer; holding out
time blocks tests temporal transfer. None establishes validity under a changed
network, timetable, fare, or crowding regime. Here grouped holdout is a model-
validity check, not a claim of forecasting performance.
